\chapter{Research Question}
\section{Problem Statement}
The problem is cross-modal alignment in a resource-constrained environment.
Given the same pretrained vision--language model, we analyse the effect of different training objectives on answer accuracy, explanation quality, and evidence sensitivity.
The model, the data format, the optimiser, and the evaluation process are kept constant wherever possible, leaving only the training signal as the variable.
The BLIP-2 pair reuses one Q-Former training setup. The Qwen2-VL generative and explanation-aware runs reuse one QLoRA setup.
The goal is not to find the best scoring setup, but to observe whether the chosen objective leads the model in a direction that matters for a more trustworthy reasoning,
and to check whether the answer quality, the explanation quality, and the evidence sensitivity move together or separately under each objective.

\section{Research Questions}
The main question is to investigate how the resource-constrained vision--language model should be aligned to achieve both accurate and faithful visual reasoning.
The following research questions break that down into more measurable parts.

\begin{enumerate}
    \item RQ1: How have vision--language alignment methods evolved for grounded visual reasoning and explanation generation?
    \item RQ2: How does BLIP-2 generative fine-tuning compare with BLIP-2 contrastive-enhanced fine-tuning on visual reasoning?
    \item RQ3: Does explanation-aware fine-tuning improve answer accuracy and explanation quality on datasets with gold rationales?
    \item RQ4: How does the same fine-tuning framework perform across document, chart, science, commonsense, and natural-image domains?
    \item RQ5: How sensitive are the model's answers to visual evidence under image-masking perturbations?
    \item RQ6: How does the selected explanation-aware QLoRA setup scale from Qwen2-VL-2B to Qwen2-VL-7B?
    \item RQ7: How well do fine-tuned adapters transfer across visual domains and prompt formats?
\end{enumerate}
